\documentclass[acmsmall,nonacm,screen]{acmart}

\begin{document}

\title{``Francis2'': Verified Algorithms in Algebraic Geometry}

\author{Hugh Coleman}
\email{htcolema@uwaterloo.ca}
\affiliation{%
  \institution{University of Waterloo}
  \state{Ontario}
  \country{Canada}
}

\renewcommand{\shortauthors}{Coleman}

% Let's spell the man's name right!
\newcommand{\Groebner}{Gr{\"o}bner}

\begin{abstract}
This project aims to implement a verified library for computational algebraic geometry within the Lean 4 proof assistant.
We plan to formalize {\it Buchberger's algorithm}, providing a certifiable mechanism for the construction of {\it \Groebner{} bases} for ideals of polynomial rings.
Using this, we will be able to efficiently decide questions of ideal membership and of radicality.

Beyond this, we will also aim to formalize the computation of the Betti numbers, a high-level homological invariant, and to provide tools for identifying Cohen--Macaulay and Gorenstein rings.
This will provide a rigorous alternative to existing unverified computer algebra systems and ad-hoc scripts.

\

\noindent\it\color{red}
This proposal was partially drafted with the assistance of a large language model, but has been reviewed and heavily edited since---turns out, they do not know much about \Groebner{} bases!
\end{abstract}

\maketitle

\section{Introduction}

Today, mathematicians rely heavily on specialized software packages to perform complex symbolic computations.
These systems are essential for tasks ranging from computing (co)homology groups to checking ideal membership.
However, there is a sharp divide between systems designed for efficient calculation and those designed for rigour: while proof assistants---such as Lean---excel at verifying abstract properties, they often lack the algorithmic machinery to perform the computations they describe.
In contrast, traditional computer algebra systems rarely offer formal guarantees that their outputs are correct.

This project seeks to leverage Lean's interactive nature to bridge this gap.
Specifically, our goal is to develop a library for performing elementary computations in commutative algebra that simultaneously produces machine-checkable proofs of its claims.
This will enable algebraists to perform analysis of algebraic varieties entirely in Lean.

\section{Background}

In his 1965 PhD thesis \cite{Buchberger1965}, Bruno Buchberger introduced the concept of a {\it \Groebner{} basis} and described an algorithm for their construction.
In effect, they are generating sets that behave predictably under multivariate polynomial division.

Informally, a \Groebner{} basis can be thought of as playing the same role for an ideal that a linear basis plays for a vector space.
Consequently, they are of interest because they provide a means for efficiently studying {\it algebraic varieties}---the ``shapes'' defined by the solution sets of polynomial equations.
For example, consider the ideal
$$
  I = (xy + y^2, xz + yz)
$$
of the ring $\mathbb{R}[x, y, z]$, which, by Buchberger's theorem, is in fact already a \Groebner{} basis for itself.
This basis induces the {\it initial ideal} $\mathrm{ini}(I) = (x y, x z)$, which can easily be seen to vanish on the $yz$-plane ($x = 0$) and on the $x$-axis ($y = 0$, $z = 0$).
In the language of geometry, this is {\bf not} {\it equidimensional} because it consists of components of different dimensions (a two-dimensional plane and a one-dimensional line)---equivalently, in the language of commutative algebra, the quotient ring $R/\mathrm{ini}(I)$ is {\bf not} {\it Cohen-Macaulay}.
Crucially, this property lifts to $R/I$, and thus we may conclude that it too must not be Cohen-Macaulay.

The Cohen-Macaulay property is a hallmark of ``regularity'', in the sense that it ensures that a variety is broadly well-behaved.
Properties like being reduced or Gorenstein, and topological invariants like the Betti numbers, capture further properties:

\begin{itemize}
\item
{\it Reducedness} ensures that the variety admits no ``infinitesmals''. For example, whereas the ideals $(x)$ and $(x^2)$ both vanish at $x = 0$, the latter admits a "fuzzy" point at the origin, in the sense that an $\varepsilon$ so small that $\varepsilon^2 = 0$ can be considered to be a vanishing point of $(x^2)$.
\item
A {\it Gorenstein} ring is a member of a family smaller than the Cohen-Macaulay rings, and captures a deep algebraic property.
\item
The {\it Betti numbers} are a topological invariant, quantifying the complexity of the relations---or {\it syzygies}---between the generators of an ideal.
\end{itemize}

\subsection{Shortcomings of Existing Software Packages}

Currently, there is no single ``gold standard'' software package for answering computational problems in commutative algebra.
While many---such as those at the University of Waterloo---opt for {\it Macaulay2}\footnote{Named for Francis Sowerby Macaulay (1862--1937), who also gives his name to this proposed project.} \cite{Macaulay2}, other specialized options such as {\it Singular}, {\it Magma}, and {\it CoCoA} coexist alongside general-purpose systems like {\it Wolfram Mathematica} and {\it SageMath}.

Despite their performance, these packages can broadly be criticized for being ``black boxes''.
Their underlying kernels are massive, unverified (or closed-source), and not invulnerable to classical software issues such as logical or runtime errors.
% Furthermore, these tools operate external to formal reasoning; one cannot easily ``lift'' a \Groebner{} basis or Betti number from Macaulay2 into a formal Lean proof. 
% This project aims to bridge this gap by providing a framework that computes these algebraic data alongside formal certificates, in an efficient and sound fashion.

\subsection{Active Work}

In September 2024, Antoine Chambert--Loir formalized monomial orders and multivariate polynomial division in Mathlib 4.
Recently, Junyu Guo and Hao Shen have been independently formalizing \Groebner{} basis theory in Lean \cite{GroebnerLean}.
Depending on time constraints, we may opt to work on top of their efforts rather than parallel to their efforts; details are in the Roadmap.

\section{Project Functionality}

% Tersely, the objective is to formalize the headline theorems and algorithms of CO 439, {\it Topics in Combinatorics: Combinatorial Commutative Algebra}.

Specifically, we intend to provide definitions of elementary structures in commutative algebra, and in particular, the definition of a \Groebner{} basis.
We will also implement proofs of Buchberger's theorem and of Hilbert's syzygy theorem; alongside algorithms for constructing \Groebner{} bases, minimal free resolutions, and Betti tables of ideals of polynomial rings.
Finally, we will build tactics capable of closing certain classes of goals---such as problems of ideal membership and of radicality---by leveraging the theory of \Groebner{} bases.

\section{Roadmap}

We plan to spend approximately three weeks on each of the following tasks.

\begin{itemize}
\item
Formalizing elementary definitions, such as monomial orderings, multivariate division, and \Groebner{} bases.
We will follow \cite{EneHerzog} (primarily for the reason that it was the textbook followed in CO 439, {\it Combinatorial Algebraic Geometry}), but will also refer to \cite{CoxLittleOShea} whenever appropriate.
\item 
Implementing Buchberger's algorithm, alongside proofs of correctness and of termination.
(The latter depends on {\it Dickson's lemma}, which we do not plan to prove ourselves.)
\item 
Implementing the computation of syzygies and the construction of minimal free resolutions and Betti tables.
Broadly, the algorithm is a straightforward induction.
\item 
Implementing high-level tactics for automatically closing goals of ideal membership and radicality.
These are tractable because they can be efficiently checked on a \Groebner{} basis.
\end{itemize}

To maximize the potential scope of this project, we intend to build upon the foundations established by Chambert-Loir.
We may also build on the proofs of Guo and Shen \cite{GroebnerLean} so that we may focus our efforts on the higher-level algorithms for syzygies, free resolutions, and Betti tables; and of the development of suitable tactics.

We anticipate that large language models will accelerate the formalization of routine lemmas. 

\bibliographystyle{ACM-Reference-Format}
\bibliography{base}

\end{document}